\section{Mục tiêu và phương pháp nghiên cứu}
\subsection{Mục tiêu nghiên cứu}
Mục tiêu của đồ án này nghiên cứu, thiết kế cơ chế tổ chức và điều khiển đội hình đa robot phối hợp di chuyển vật thể cần thao theo quỹ đạo mong muốn thông qua tương tác đẩy, trong đó bao gồm việc xây dựng cấu trúc đội hình và việc thiết kế chiến thuật đẩy. Bên cạnh đó, đồ án cũng đề xuất một phương pháp tiếp cận mới để ước lượng vị trí của điểm trọng tâm. Các mục tiêu cụ thể như sau

\begin{itemize}
    \item Nghiên cứu cơ chế xây dựng và chuyển đổi đội hình đa robot cho mục tiêu đẩy vật dựa trên các yếu tố về động lực học, nhằm tạo ra được một cấu trúc ổn định nhất trong quá trình vận chuyển, đảm bảo sự linh hoạt và khả năng thích ứng cao khi có sự thay đổi về số lượng robot tham gia hoặc wrench mong muốn.
    \item Nghiên cứu chiến thuật đẩy vật thể của hệ đa robot tương ứng với đội hình ở thời điểm hiện tại, sao cho tổng hợp lực-mô men sinh ra bởi toàn hệ robot đạt được wrench mong muốn.
    \item Đề xuất một phương pháp ước lượng vị trí điểm trọng tâm bằng cách đẩy thử dựa trên yếu tố về gia tốc góc.
\end{itemize}

\subsection{Phạm vi nghiên cứu}
Trong nghiên cứu này, yếu tố về hình dạng của vật thể cần thao tác sẽ được giả định là đã biết trước sau quá trình thu thập và phân tích dữ liệu từ các cảm biến ngoại vi, giới hạn trong phạm vi các vật thể có dạng đa giác (polygon). Vị trí và hướng của robot được tính toán trực tiếp mà không sử dụng các bộ lọc như trong các hệ thống thực tế, robot cũng được giả định có khả năng duy trì tương tác với vật thể trong suốt thời gian đẩy, giới hạn chỉ được đẩy theo hướng vuông góc với cạnh tiếp xúc. Hệ thống robot là phân tán, không có bộ xử lý trung tâm, tuy nhiên vẫn cho phép các robot truyền thông với nhau. Việc di chuyển đến vị trí tương của các robot giữa mỗi lần chuyển đổi đội hình được bỏ qua, coi như có thể di chuyển đến vị trí mong muốn mà không xảy ra va chạm.

Các giải pháp được đề xuất trong đồ án sẽ được đánh giá và thử nghiệm thông qua mô phỏng máy tính hoặc một số mô hình thực nghiệm giả định. Hệ thống không được thử nghiệm trực tiếp trên các robot thực tế, mà chỉ tập trung vào các điều kiện giả lập nhằm kiểm tra tính khả thi và hiệu quả của các thuật toán đề xuất.
    
\subsection{Phương pháp nghiên cứu}

Đồ án sử dụng ngôn ngữ lập trình Python để mô phỏng hệ thống đa robot trong môi trường 2D, ma sát giữa các bề mặt được mô hình hoá bằng mô hình ma sát Coulomb, trong đó lực ma sát tỉ lệ với lực pháp tuyến và có hướng ngược với chuyển động tương đối.

Thuật toán xây dựng đội hình và lên chiến thuật đẩy được dựa trên nghiên cứu \cite{Rosenfelder2024}, trong đó sự ổn định của đội hình được đánh giá trên không gian wrench space dựa trên lý thuyết về nắm bắt (grasping) bằng một phương pháp lấy cảm hứng từ \cite{ferrari1992planning} và được giải cục bộ bằng thuật toán tối ưu bầy đàn hạt dùng Lagrangian tăng cường trong môi trường phân tán (Distributed Augmented Lagrangian Particle Swarm Optimization - DALPSO), còn lực đẩy của mỗi robot được tính toán cục bộ thông qua việc giải một bài toán tối ưu hàm bậc hai với ràng buộc tuyến tính (Quadratic problem - QP).

Phương pháp ước lượng vị trí trọng tâm được đề xuất dựa trên việc đo lường gia tốc góc của vật thể thông qua các lần đẩy thử, từ đó xây dựng hệ phương trình dựa trên mô hình ma sát Coulomb và giải bằng phương pháp bình phương tối thiểu (Least Squares - LS).